\clearpage\thispagestyle{empty}
\pdfbookmark[0]{Sources for the Letter to Bode}{bode18060730-sources}
\begin{center}\textbf{Sources for the Letter to Bode}\end{center}
\noindent This AI translation follows Carl Friedrich Gauss’s letter of 30 July 1806 to Johann Elert Bode as printed in \href{https://gdz.sub.uni-goettingen.de/id/PPN236007467}{Gauss’s \textit{Werke}, vol.~VI} (Göttingen, 1874), pp.~278--279. The letter appeared in Bode’s \textit{Astronomisches Jahrbuch für 1809} (Berlin, 1806), \href{https://www.digitale-sammlungen.de/en/view/bsb10538338?page=225}{pp.~215}--\href{https://www.digitale-sammlungen.de/en/view/bsb10538338?page=229}{219}. The two historical notes also appear on \href{https://www.digitale-sammlungen.de/en/view/bsb10538338?page=227}{p.~217 of the yearbook}. The second is signed by Bode.
\par\medskip
The yearbook prints the Pallas ephemeris on \href{https://www.digitale-sammlungen.de/en/view/bsb10538338?page=228}{p.~218} and the Juno ephemeris on \href{https://www.digitale-sammlungen.de/en/view/bsb10538338?page=229}{p.~219}. The Juno table there begins on 12 January 1807. The bracketed description in the collected edition gives 7 December 1806.
\par\medskip
The Milan observations of Juno and Pallas are in the July 1806 \textit{Monatliche Correspondenz}, vol.~XIV, on \href{https://www.digitale-sammlungen.de/en/view/bsb10538606?page=98}{p.~88} and \href{https://www.digitale-sammlungen.de/en/view/bsb10538606?page=100}{p.~90}, respectively. The journal credits Carlini with making and computing the Pallas observations and the Juno observations at the equatorial sector. The historical scans are from the Bayerische Staatsbibliothek. Translation and typesetting produced by AI.
\clearpage
